%%%%%
%
\section{Polynome de premier degr\'e dans un domaine}

\itemstyle{Op\'erateur~:} \operfmt{Poly\_1\_flt\_domain}
\index{Poly\_1\_flt\_domain}

%%%%%OBJET%%%%
\itemstyle{Objet~:}
Application d'un polyn\^ome de 1$^{er}$ degr\'e
\`a un parametre sur un domaine (seulement si le param\`etre appartient \`a un domaine).

%%%%%ARGUMENTS%%%%
\bigskip
%%%%%ENTREES%%%%
\begin{tabarg}
\multicolumn{3}{l}{\titretab {Entr\'ees}}\\
\hline
$X$ & Vecteur, \texttt{float} & Op\'erande (unit\'e quelconque)\\
$C_0$ & Coeff., \texttt{float} & Coefficient du terme de degr\'e 0.\\
$C_1$ & Coeff., \texttt{float} & Coefficient du terme de degr\'e 1.\\
$Borne\_{min}$ & Coeff., \texttt{float} & Borne minimale pour le param\`etre X\\
$Borne\_{max}$ & Coeff., \texttt{float} & Borne maximale pour le param\`etre X\\
$Unit\_{res}$ & Coeff., \texttt{String} & Unit\'e du r\'esultat.\\
$Name $ & Coeff., \texttt{string} & Descriptif du r\'esultat\\
$Eng\_name$ & Coeff., \texttt{string} & Descriptif du r\'esultat en langue anglaise\\
\end{tabarg}

\medskip
%%%%%SORTIES%%%%
\begin{tabarg}
\multicolumn{3}{l}{\titretab {Sorties}}\\
\hline
$Z$ & Vecteur, \texttt{float} & R\'esultat (unit\'e indiqu\'ee par $Unit\_{res}$)\\
\end{tabarg}


%%%%%FORMULE%%%%
\itemstyle{Formule}

%% Poly_1_flt_domain

\[
Z = 
\begin{cases}
C_0 + C_1. X \qquad \text{ si } Borne_{min} \le X \le Borne_{max} \\
\text{Une valeur ind�finie al�atoire} \qquad \text{ sinon }  \\
\end{cases}
\]



%%%%%COMMENTAIRE%%%%
\itemstyle{Commentaire~:}
On effectue le calcul uniquement si la variable X est dans le domaine de d\'efinition :
 c'est \`a dire sup\'erieure ou \'egale \`a Borne\_mini et inf\'erieure ou \'egale \`a Borne\_maxi.

%%%%%AUTEUR%%%%
\itemstyle{Auteur~:}
CNRM/GMEI/TRAMM

%%%%%REFERENCE%%%%
\itemstyle{R\'ef\'erences~:}
